% Part: second-order-logic
% Chapter: metatheory
% Section: compactness

\documentclass[../../../include/open-logic-section]{subfiles}

\begin{document}

\olfileid{sol}{met}{com}

\olsection{દ્વિતીય-ક્રમ તર્કશાસ્ત્ર સઘન નથી}

\begin{explain}
વાક્યોનો ગણ~$\Gamma$ \emph{સાન્ત રીતે સંતોષ્ય} ત્યારે કહેવાય જ્યારે
તેનો દરેક સાન્ત ઉપગણ સંતોષ્ય હોય. પ્રથમ-ક્રમ તર્કશાસ્ત્રમાં ગુણધર્મ છે
કે !!{sentence}sનો ગણ~$\Gamma$ સાન્ત રીતે સંતોષ્ય હોય તો તે પોતે
સંતોષ્ય છે. આ ગુણધર્મને \emph{સઘનતા} કહે છે. ફલિતતા દ્વારા તેનું
સમકક્ષ સ્વરૂપ છે: $\Gamma \Entails !A$ હોય તો કોઈ સાન્ત ઉપગણ
$\Gamma_0 \subseteq \Gamma$ માટે પણ $\Gamma_0 \Entails
!A$. આ સ્વરૂપમાં તે પૂર્ણતા પ્રમેયનું તત્કાલ અનુસિદ્ધાંત છે: જો
$\Gamma \Entails !A$, તો પૂર્ણતા મુજબ $\Gamma \Proves !A$. પરંતુ
!!{derivation},~$\Gamma$નાં માત્ર સાન્ત સંખ્યામાં !!{sentence}sનો
ઉપયોગ કરી શકે છે.

દ્વિતીય-ક્રમ તર્કશાસ્ત્રમાં સઘનતા સાચી નથી. દ્વિતીય-ક્રમ
!!{sentence}sના એવા ગણો છે જે સાન્ત રીતે સંતોષ્ય છે પરંતુ પોતે
સંતોષ્ય નથી; અને એવા ગણો પણ છે જે અમુક~$!A$ને ફલિત કરે છે,
પરંતુ તેમનો કોઈ સાન્ત ઉપગણ~$!A$ને ફલિત કરતો નથી.
\end{explain}


\begin{thm}
\ollabel{thm:sol-not-compact}
દ્વિતીય-ક્રમ તર્કશાસ્ત્ર સઘન નથી.
\end{thm}
\sourcecorrection{OLSOL-011}{સ્થિર અંગ્રેજી મૂળમાં આ પ્રમેયનું
\texttt{thm:sol-undecidable} લેબલ અગાઉના અનિર્ણેયતા પ્રમેયનું લેબલ
જ ફરી વાપરે છે. બે અલગ પરિણામોને એક જ લેબલ ન રહે તે માટે અહીં
\texttt{thm:sol-not-compact} આપ્યું છે.}

\begin{proof}
યાદ કરો કે
\[
\fn{Inf} \ident \lexists[u][(\lforall[x][\lforall[y][(\eq[u(x)][u(y)]
      \lif \eq[x][y])]] \land \lexists[y][\lforall[x][\eq/[y][u(x)]]])]
\]
એવું વાક્ય છે જે !!a{structure}માં ત્યારે અને માત્ર ત્યારે જ
સંતોષાય જ્યારે તેનું વ્યાપકક્ષેત્ર અનંત હોય. $!A^{\ge n}$ એવું
વાક્ય હોય જે વ્યાપકક્ષેત્રમાં ઓછામાં ઓછા $n$ !!{element}s છે એમ કહે;
દાખલા તરીકે,
\[
!A^{\ge n} \ident \lexists[x_1][\dots\lexists[x_n][(\eq/[x_1][x_2]
    \land \eq/[x_1][x_3] \land \dots \land \eq/[x_{n-1}][x_n])]].
\]
આ !!{sentence}sનો ગણ વિચારો:
\[
\Gamma = \{\lnot \fn{Inf}, !A^{\ge 1}, !A^{\ge 2}, !A^{\ge 3}, \dots\}.
\]
તે સાન્ત રીતે સંતોષ્ય છે: મનસ્વી સાન્ત ઉપગણ~$\Gamma_0
\subseteq \Gamma$ માટે ઓછામાં ઓછો એક $k$ એવો પસંદ કરી શકાય કે
$!A^{\ge k} \in \Gamma$ અને $n>k$ માટે પસંદ કરેલા સાન્ત ઉપગણમાં
કોઈ $!A^{\ge n}
\in \Gamma_0$ ન હોય. જો $\Domain{M}$માં $k$ !!{element}s હોય,
તો $\Sat{M}{\Gamma_0}$. પરંતુ~$\Gamma$ સંતોષ્ય નથી: જો
$\Sat{M}{\lnot \fn{Inf}}$, તો~$\Domain{M}$ સાન્ત હોવું જોઈએ;
માનો તેનું કદ~$k$ છે. ત્યારે $\Sat/{M}{!A^{\ge k+1}}$.
\sourcecorrection{OLSOL-012}{સ્થિર અંગ્રેજી મૂળ સાન્ત ઉપગણની
દલીલમાં $n>k$ માટે $!A^{\ge n}$ આખા $\Gamma$માં નથી એમ કહે છે,
પરંતુ એ બધા વાક્યો $\Gamma$માં છે. અહીં તે શરત પસંદ કરેલા સાન્ત
ઉપગણ~$\Gamma_0$ માટે મૂકી છે. તેમાં કદની એક પણ શરત ન હોય તો
ઓછામાં ઓછું એક ઘટક ધરાવતું નિદર્શ લેવા $k$ને એક લઈ શકાય છે.}
\end{proof}

\begin{prob}
એવો ગણ~$\Gamma$ અને !!a{sentence}~$!A$નું ઉદાહરણ આપો કે
$\Gamma \Entails !A$, પરંતુ દરેક સાન્ત ઉપગણ~$\Gamma_0 \subseteq
\Gamma$ માટે $\Gamma_0 \Entails/ !A$.
\end{prob}


\end{document}
